% DOCUMENT CLASS SETUP
\documentclass[12pt, a4paper]{article}

% DOCUMENT PACKAGES
\usepackage[a4paper, margin=1in]{geometry}
\usepackage{titling}
\usepackage{enumitem}
\usepackage{amsmath}
\usepackage{graphicx}
\usepackage{xcolor}
\usepackage{listings}
\usepackage{hyperref}

% FONT AND ENCODING SETUP
\usepackage[english]{babel}
\usepackage[utf8]{inputenc}
\usepackage[T1]{fontenc}
\usepackage{lmodern}
\usepackage{fontspec}

% PROFESSIONAL CODE BLOCK
\usepackage{tcolorbox}
\tcbuselibrary{minted, skins, breakable}

% DEFINE GITHUB DARK THEME COLORS %
\definecolor{ghback}{HTML}{0d1117}    % Dark background
\definecolor{ghtext}{HTML}{c9d1d9}    % Light gray text
\definecolor{ghborder}{HTML}{30363d}  % GitHub Border color

% PROFESSIONAL CODE BLOCK CONFIGURATION (MINTED VERSION) %
\renewcommand{\theFancyVerbLine}{\textcolor{white!60!gray}{\scriptsize\arabic{FancyVerbLine}}}
\newtcblisting{codeblock}[2][]{
    enhanced,
    breakable,
    colback=ghback,
    colframe=ghborder,
    coltitle=ghtext,
    colupper=ghtext,
    fonttitle=\bfseries\small,
    arc=3mm,
    boxrule=0.5mm,
    % PADDING SETTINGS
    left=25pt,
    right=10pt,
    top=5pt,
    bottom=5pt,
    % MINTED SETTINGS
    listing engine=minted,
    listing only,
    minted language=python,
    minted options={
        linenos,
        numbersep=5pt,
        fontsize=\footnotesize,
        breaklines=true,
        breakanywhere=true,
        autogobble,
        style=monokai,
    },
    title={#2},
    #1
}

% DOCUMENT TITLE SETUP %
\title {
	\textbf{2110452 HIGH PERFORMANCE ARCHITECTURE} \\
    \vspace{1em}
	\Large{\textbf{Homework 1: Vectorization}} \\
    \vspace{0.5em}
    \large{Computer Engineering, Chulalongkorn Univerity}
}
\author{Worralop Srichainont}
\date{\today}
\setlength{\droptitle}{15em}

% DOCUMENT CONTENTS
\begin{document}

    \maketitle
    \thispagestyle{empty}
    \clearpage

    \tableofcontents
    \thispagestyle{empty}
    \clearpage

    \pagenumbering{arabic}

    \section{Problem I}
        For this problem, we are running the experiment on two \texttt{C} programs.
        \begin{itemize}
            \item \texttt{no\_avx.c} is the \texttt{C} program to add 2 integer arrays without vectorization.
            \item \texttt{avx.c} is the \texttt{C} program to add 2 integer arrays with vectorization.
        \end{itemize}

        \subsection{No Vectorization}
            For \texttt{C} libraries, we are going to use.
            \begin{itemize}
                \item \texttt{stdio.h} for basic I/O commands.
                \item \texttt{stdlib.h} for memory allocation using \texttt{malloc()} and \texttt{free()}.
                \item \texttt{time.h} to set a timer for an experiment.
            \end{itemize}

            \begin{codeblock}[minted language=c]{Libraries}
                #include<stdio.h>
                #include<stdlib.h>
                #include<time.h>
            \end{codeblock}

            First, we are define some constants for this test.
            \begin{itemize}
                \item \texttt{TEST\_ITERATIONS} is amount of times we execute the experiment.
                \item \texttt{PROCESS\_ITERATIONS} is amount of times we add 2 arrays for significant runtime differences.
                \item \texttt{ARRAY\_SIZE} is length of the integer array.
            \end{itemize}

            \begin{codeblock}[minted language=c]{Constants}
                #define TEST_ITERATIONS 10
                #define PROCESS_ITERATIONS 100
                #define ARRAY_SIZE 6400000
            \end{codeblock}

            Then, write a simple arrays addition program as follows:
            \begin{codeblock}[minted language=c]{Scalar Addition}
                void add(int *a, int *b) {
                    for(int i = 0; i < ARRAY_SIZE; i++) {
                        a[i] += b[i];
                    }
                }
            \end{codeblock}

    \pagebreak

            Then, write a \texttt{run\_test()} function to calculate the time for non-vectorized arrays addition.
            \begin{itemize}
                \item Use \texttt{malloc()} to allocate a single block of memory for 6,400,000 integers.
                \item Initialize values inside of both integer arrays.
                \item Run an experiment for 100 times to see the significant differences.
                \item Free the allocated memory.
                \item Convert \texttt{clock\_t} value into \texttt{double} type in second unit, and return it.
            \end{itemize}

            \begin{codeblock}[minted language=c]{Runtime Calculation}
                double run_test() {
                    // Allocated memory for array a and b.
                    int *a = (int*)malloc(ARRAY_SIZE * sizeof(int));
                    int *b = (int*)malloc(ARRAY_SIZE * sizeof(int));

                    // Initialize value inside each array.
                    for(int i = 0; i < ARRAY_SIZE; i++) {
                        a[i] = i;
                        b[i] = i;
                    }

                    // Run an experiment
                    clock_t begin_time = clock();
                    for(int i = 0; i < PROCESS_ITERATIONS; i++) {
                        add(a, b);
                    }
                    clock_t end_time = clock();

                    // Free allocated memory
                    free(a);
                    free(b);

                    // Return executed time
                    double run_time = ((double) end_time - begin_time) / CLOCKS_PER_SEC;
                    return run_time;
                }
            \end{codeblock}

            Finally, write a \texttt{main()} function to loop the experiment 10 times and display executed time for each one.
            \begin{codeblock}[minted language=c]{Run an Experiment}
                int main() {
                    printf("============ SCALAR RUNTIME ===========\n");
                    for(int i = 1; i <= TEST_ITERATIONS; i++) {
                        printf("TEST %d: %f s\n", i, run_test());
                    }
                    return 0;
                }
            \end{codeblock}

    \pagebreak

        \subsection{Vectorization}
            For \texttt{C} libraries, we are going to use.
            \begin{itemize}
                \item \texttt{stdio.h} for basic I/O commands.
                \item \texttt{stdlib.h} for memory allocation using \texttt{malloc()} and \texttt{free()}.
                \item \texttt{immintrin.h} allows direct use of CPU instructions without writing Assembly.
                \item \texttt{time.h} to set a timer for an experiment.
            \end{itemize}

            \begin{codeblock}[minted language=c]{Libraries}
                #include<stdio.h>
                #include<stdlib.h>
                #include<immintrin.h>
                #include<time.h>
            \end{codeblock}

            First, we are define some constants for this test.
            \begin{itemize}
                \item \texttt{TEST\_ITERATIONS} is amount of times we execute the experiment.
                \item \texttt{PROCESS\_ITERATIONS} is amount of times we add 2 arrays for significant runtime differences.
                \item \texttt{ARRAY\_SIZE} is length of the integer array.
            \end{itemize}

            \begin{codeblock}[minted language=c]{Constants}
                #define TEST_ITERATIONS 10
                #define PROCESS_ITERATIONS 100
                #define ARRAY_SIZE 6400000
            \end{codeblock}

            For the next step, please read the next page $\rightarrow$

    \pagebreak

            Then, write a vectorized version of array addition by using loop unroll technique.
            \begin{enumerate}
                \item Create 256-bit register (\texttt{\_\_m256i}) to store 8 integers (32 bits each).
                \item The function \texttt{\_mm256\_loadu\_si256} grabs 8 integers (256 bits) from data memory at once and puts them into the CPU vector register.
                \item The function \texttt{\_mm256\_add\_epi32} is the parallel add function to add 8 integers simultaneously in a single clock cycle.
                \item The function \texttt{\_mm256\_storeu\_si256} write the sum of 8 integers back into an array.
                \item After done addition, go to the next 8 integers.
            \end{enumerate}

            \begin{codeblock}[minted language=c]{Vectorized Addition Part I}
                void add_avx(int *a, int *b) {
                    int i = 0;
                    for(; i < ARRAY_SIZE; i += 8) {
                        // Load 256 bits chunks of each array.
                        __m256i av = _mm256_loadu_si256((__m256i*) &a[i]);
                        __m256i bv = _mm256_loadu_si256((__m256i*) &b[i]);

                        // Add each pair of 32 bits integers in chunks.
                        av = _mm256_add_epi32(av, bv);

                        // Store 256 bits chunks to variable a.
                        _mm256_storeu_si256((__m256i*) &a[i], av);
                    }
                }
            \end{codeblock}

            \begin{enumerate}
                \item[6.] For the leftovers (e.g. array of 26 integers would have 2 leftover integers), just use non-vectorized addition.
            \end{enumerate}

            \begin{codeblock}[minted language=c]{Vectorized Addition Part II}
                void add_avx(int *a, int *b) {
                    /* INSERT PART I CODE... */
                    for(; i < ARRAY_SIZE; i++) {
                        a[i] += b[i];
                    }
                }
            \end{codeblock}

    \pagebreak

            Then, write a \texttt{run\_test\_avx()} function to calculate the time for vectorized arrays addition.
            \begin{itemize}
                \item Use \texttt{malloc()} to allocate a single block of memory for 6,400,000 integers.
                \item Initialize values inside of both integer arrays.
                \item Run an experiment for 100 times to see the significant differences.
                \item Free the allocated memory.
                \item Convert \texttt{clock\_t} value into \texttt{double} in second unit, and return it.
            \end{itemize}

            \begin{codeblock}[minted language=c]{Runtime Calculation}
                double run_test_avx() {
                    // Allocated memory for array a and b.
                    int *a = (int*)malloc(ARRAY_SIZE * sizeof(int));
                    int *b = (int*)malloc(ARRAY_SIZE * sizeof(int));

                    // Initialize value inside each array.
                    for(int i = 0; i < ARRAY_SIZE; i++) {
                        a[i] = i;
                        b[i] = i;
                    }

                    // Run experiment
                    clock_t begin_time = clock();
                    for(int i = 0; i < PROCESS_ITERATIONS; i++) {
                        add_avx(a, b);
                    }
                    clock_t end_time = clock();

                    // Free allocated memory
                    free(a);
                    free(b);

                    // Return executed time
                    double run_time = ((double) end_time - begin_time) / CLOCKS_PER_SEC;
                    return run_time;
                }
            \end{codeblock}

            Finally, write a \texttt{main()} function to loop the experiment 10 times and display executed time for each one.
            \begin{codeblock}[minted language=c]{Run an Experiment}
                int main() {
                    printf("============ VECTORIZE RUNTIME ===========\n");
                    for(int i = 1; i <= TEST_ITERATIONS; i++) {
                        printf("TEST %d: %f s\n", i, run_test_avx());
                    }
                    return 0;
                }
            \end{codeblock}

        \subsection{Run Experiment}

            By running both \texttt{C} files at the same time, we can write a batch script.
            \begin{itemize}
                \item We are using GCC compiler to compile the \texttt{C} program into \texttt{.exe} file.
                \item Write \texttt{gcc -O3} to indicate optimization level of the GCC compiler.
                \item For \texttt{no\_avx.c} file, use \texttt{-fno-tree-vectorize} to prevent compiler from performing vectorization.
                \item For \texttt{avx.c} file, use \texttt{-mavx2} to allows using \texttt{immintrin.h} library.
            \end{itemize}
            \begin{codeblock}[minted language=bat]{Batch Script}
                REM COMPILE PROGRAM USING GCC COMPILER
                gcc -O3 -fno-tree-vectorize no_avx.c -o no_avx.exe
                gcc -O3 -mavx2 avx.c -o avx.exe

                REM EXECUTE EXPERIMENT
                no_avx.exe
                avx.exe

                PAUSE
            \end{codeblock}

            After that, just run this batch script and everything is done!

    \pagebreak

        \subsection{Result}

            After the running the batch script, the program shows the following result.

            \begin{table}[h]
                \centering
                \begin{tabular}{|c|c|c|}
                    \hline
                    \rule{0pt}{2.5ex}
                    \textbf{No.} & \textbf{No Vectorization} & \textbf{Vectorization} \\
                    \hline \hline
                    1 & 0.111 s & 0.015 s \\
                    2 & 0.133 s & 0.019 s \\
                    3 & 0.131 s & 0.017 s \\
                    4 & 0.120 s & 0.016 s \\
                    5 & 0.117 s & 0.016 s \\
                    6 & 0.126 s & 0.016 s \\
                    7 & 0.114 s & 0.015 s \\
                    8 & 0.114 s & 0.016 s \\
                    9 & 0.115 s & 0.014 s \\
                    10 & 0.112 s & 0.017 s \\
                    \hline \hline
                    \textbf{Trimmed mean} & 0.119 s & 0.016 s \\
                    \hline
                \end{tabular}
                \caption{Experiment Result}
            \end{table}

            From the trimmed mean of both version, we can easily calculate the speedup of the vectorization as follows:
            \begin{align*}
                \text{Speedup} &= \frac{{\text{Time}}_{\text{Scalar}}}{{\text{Time}}_{\text{AVX}}} \\[1em]
                \text{Speedup} &= \frac{0.119 \ \text{s}}{0.016 \ \text{s}} \\[1em]
                \text{Speedup} &\approx 7.44
            \end{align*}

            The experiment shows that the vectorization for array addition in \texttt{C} language is \textbf{7.44} times faster.

    \pagebreak

    \section{Problem II}
        For this problem, we are running the experiment on two \texttt{python} programs.
        \begin{itemize}
            \item \texttt{no\_avx.py} is the \texttt{python} program to add 2 integer arrays without vectorization.
            \item \texttt{avx.py} is the \texttt{python} program to add 2 integer arrays with vectorization.
        \end{itemize}

        \subsection{No Vectorization}
            For \texttt{python} library, we need to use \texttt{time} library for to set a timer for an experiment.
            \begin{codeblock}{Libraries}
                import time
            \end{codeblock}

            First, we are define some constants for this test.
            \begin{itemize}
                \item \texttt{TEST\_ITERATIONS} is amount of times we execute the experiment.
                \item \texttt{ARRAY\_SIZE} is length of the integer array.
            \end{itemize}

            \begin{codeblock}[minted language=c]{Constants}
                TEST_ITERATIONS = 10
                ARRAY_SIZE = 6400000
            \end{codeblock}

            Then, write a \texttt{run\_test()} function to calculate the time for non-vectorized arrays addition.
            \begin{itemize}
                \item Initialize two lists, each has 6,400,000 numbers.
                \item Iterate each number in the list and calculate the sum.
                \item Return the executed time.
            \end{itemize}

            \begin{codeblock}{Runtime Calculation}
                def run_test():
                    # Create two integer arrays
                    a = list(range(ARRAY_SIZE))
                    b = list(range(ARRAY_SIZE))

                    # Run an experiment
                    begin_time = time.time()
                    for i in range(ARRAY_SIZE):
                        a[i] += b[i]
                    end_time = time.time()

                    # Return executed time
                    return end_time - begin_time
            \end{codeblock}

    \pagebreak

            Finally, write a \texttt{main()} function to loop the experiment 10 times and display executed time for each one.
            \begin{codeblock}{Run an Experiment}
                def main():
                    print("============ SCALAR RUNTIME ===========")
                    for i in range(TEST_ITERATIONS):
                        print(f"TEST {i + 1}: {run_test():.3f} s")

                if __name__ == "__main__":
                    main()
            \end{codeblock}

    \pagebreak

        \subsection{Vectorization}
            For \texttt{python} libraries, we are going to use.
            \begin{itemize}
                \item \texttt{time} to set a timer for an experiment.
                \item \texttt{numpy} for vectorized mathematics operations.
            \end{itemize}

            \begin{codeblock}{Libraries}
                import time
                import numpy as np
            \end{codeblock}

            First, we are define some constants for this test.
            \begin{itemize}
                \item \texttt{TEST\_ITERATIONS} is amount of times we execute the experiment.
                \item \texttt{ARRAY\_SIZE} is length of the integer array.
            \end{itemize}

            \begin{codeblock}[minted language=c]{Constants}
                TEST_ITERATIONS = 10
                ARRAY_SIZE = 6400000
            \end{codeblock}

            Then, write a \texttt{run\_test()} function to calculate the time for vectorized arrays addition.
            \begin{itemize}
                \item Initialize two \texttt{numpy} array, each has 6,400,000 numbers.
                \item Simply add two \texttt{numpy} array. This is vectorized addition.
                \item Return the executed time.
            \end{itemize}

            \begin{codeblock}{Runtime Calculation}
                def run_test():
                    # Create two integer arrays
                    a = np.arange(0, ARRAY_SIZE)
                    b = np.arange(0, ARRAY_SIZE)

                    # Run an experiment
                    begin_time = time.time()
                    a += b
                    end_time = time.time()

                    # Return executed time
                    return end_time - begin_time
            \end{codeblock}

    \pagebreak

            Finally, write a \texttt{main()} function to loop the experiment 10 times and display executed time for each one.
            \begin{codeblock}{Run an Experiment}
                def main():
                    print("============ VECTORIZE RUNTIME ===========")
                    for i in range(TEST_ITERATIONS):
                        print(f"TEST {i + 1}: {run_test():.3f} s")

                if __name__ == "__main__":
                    main()
            \end{codeblock}

    \pagebreak

        \subsection{Run Experiment}
            By running both \texttt{python} files at the same time, we can write a batch script as follows:
            \begin{codeblock}[minted language=bat]{Batch Script}
                REM RUN EXPERIMENT
                python -u no_avx.py
                python -u avx.py

                PAUSE
            \end{codeblock}

            After that, just run this batch script and everything is done!

    \pagebreak

        \subsection{Result}
            After the running the batch script, the program shows the following result.

            \begin{table}[h]
                \centering
                \begin{tabular}{|c|c|c|}
                    \hline
                    \rule{0pt}{2.5ex}
                    \textbf{No.} & \textbf{No Vectorization} & \textbf{Vectorization} \\
                    \hline \hline
                    1 & 0.658 s & 0.008 s \\
                    2 & 0.624 s & 0.005 s \\
                    3 & 0.674 s & 0.006 s \\
                    4 & 0.688 s & 0.005 s \\
                    5 & 0.561 s & 0.007 s \\
                    6 & 0.616 s & 0.006 s \\
                    7 & 0.623 s & 0.008 s \\
                    8 & 0.600 s & 0.006 s \\
                    9 & 0.580 s & 0.005 s \\
                    10 & 0.585 s & 0.005 s \\
                    \hline \hline
                    \textbf{Trimmed mean} & 0.620 s & 0.006 s \\
                    \hline
                \end{tabular}
                \caption{Experiment Result}
            \end{table}

            From the trimmed mean of both version, we can easily calculate the speedup of the vectorization as follows:
            \begin{align*}
                \text{Speedup} &= \frac{{\text{Time}}_{\text{Scalar}}}{{\text{Time}}_{\text{AVX}}} \\[1em]
                \text{Speedup} &= \frac{0.620 \ \text{s}}{0.006 \ \text{s}} \\[1em]
                \text{Speedup} &\approx 103.33
            \end{align*}

            The experiment shows that the vectorization for array addition in \texttt{python} language is \textbf{103.33} times faster.

    \pagebreak

    \section{Problem III}

    \noindent \textbf{Question}: While vectorization is powerful, please explain a situation when it may not be beneficial.

    \vspace{1em}

    \noindent \textbf{Answer}: There are some situation where vectorization may not beneficial, such as:
    \begin{enumerate}
        \item \textbf{Data Dependencies}: The program requires the result from previous iteration to calculate, so it cannot execute in parallel. (e.g. \texttt{a[i] = a[i - 1] + b[i]})
        \item \textbf{Conditional Branching}: Vectorization applies the same instructions into multiple data simultaneously. If inside the loop has complex \texttt{if-else} statements, each element would have different action and vectorization is definitely not suitable for the task like this.
        \item \textbf{Non-Contiguous Memory}: Vectorization works best when the data are contiguous memory like \texttt{array} in \texttt{C} and \texttt{list} in \texttt{python}. If the data are scattered, vectorization would be slower since it needs to gather all the scattered data before doing the actual operation.
    \end{enumerate}

\end{document}